% Generated draft: requires mathematical and source review.
\documentclass[11pt]{article}
\usepackage[margin=25mm]{geometry}
\usepackage{fontspec}
\setmainfont{Latin Modern Roman}
\usepackage{amsmath,amssymb,mathtools,unicode-math}
\setmathfont{Latin Modern Math}
\usepackage{enumitem,needspace,xurl,hyperref}
\setlength{\parindent}{0pt}
\setlength{\parskip}{5pt}
\setlength{\emergencystretch}{3em}
\newcommand{\N}{\mathbb{N}}
\newcommand{\Z}{\mathbb{Z}}
\newcommand{\Q}{\mathbb{Q}}
\newcommand{\R}{\mathbb{R}}
\newcommand{\C}{\mathbb{C}}
\newcommand{\F}{\mathbb{F}}
\newcommand{\A}{\mathbb{A}}
\newcommand{\PP}{\mathbb{P}}
\newcommand{\E}{\mathbb{E}}
\newcommand{\eps}{\varepsilon}
\newcommand{\dd}{\,\mathrm{d}}
\begin{document}
\section*{UnsolvedMath ID: 30002409}
\textbf{Source problem code:} catalogue entry
\subsection{Definitions and mathematical statement}
\textbf{Research question.} Can noncommutative geometry provide a better understanding of automorphic spaces of cusp forms and Eisenstein series as analytic spaces, perhaps through a noncommutative boundary?

\textbf{Terminology and unresolved scope.} The expressions ``cusp forms'' and ``Eisenstein series'' designate the automorphic objects mentioned in the supplied record. The record does not specify an underlying group, base field, arithmetic subgroup, automorphic quotient, weight, level, or other automorphic data. It also does not specify whether the intended analytic spaces are spaces of automorphic functions, parameter spaces for such functions, or spaces on which they are realized. Consequently, no particular domain or analytic-space structure can be defined from the supplied information.

The proposed noncommutative boundary is a possible approach, not a required or defined construction. Neither the intended notion of noncommutative space nor the property constituting a ``better understanding'' is specified. The question therefore records a research direction; the supplied evidence does not determine a precise quantified conjecture or a construction with stated success criteria.
\subsection{Short English statement}
The question asks whether noncommutative geometry, possibly using a noncommutative boundary, can clarify the analytic interpretation of automorphic cusp forms and Eisenstein series. The underlying automorphic data, meaning of analytic space, and desired improvement remain unspecified.
\subsection{Sources}
\begin{itemize}
\item Supplied local UnsolvedMath registry record: permanent ID \(30002409\), title ``Noncommutative Geometry of Automorphic Function Spaces,'' collection ``Oberwolfach Reports Open Problems,'' exact source problem code \texttt{catalogue entry}. No individual proposer is identified; the supplied proposed year is 2013.
\item Supplied report citation: \emph{Noncommutative Geometry (2014)}, Oberwolfach Reports, DOI 10.4171/owr/2013/45, pp.\ 2553--2629. Supplied source link: \url{https://doi.org/10.4171/owr/2013/45}. The citation is preserved as supplied; its 2014 title date and the 2013 provenance require reconciliation.
\item Supplied literature source, labeled ``Oberwolfach Report'': \url{https://ems.press/content/serial-article-files/46476}.
\item Literature assessment dated exactly \texttt{2026-08-21}: the supplied curation labels the question open and describes it as a broad request for a better analytic-space interpretation, without a precise conjecture or a demonstrated general construction. This is supplied dated curation, not an independently verified literature assessment. The background expressly requests independent verification of time-sensitive status and cited claims.
\item The registry labels the statement \texttt{corrected\_verified} and says that mathematical content was unchanged while notation, typography, and wording were normalized. These are supplied editorial assertions, not independent mathematical verification of this draft.
\end{itemize}
\end{document}
